\documentclass{article}
\usepackage{lmodern}
\usepackage{amsmath}
\usepackage{listings}
\usepackage{fullpage}
\usepackage{parskip}
\usepackage{xcolor}
\lstset{
	basicstyle=\ttfamily,
	columns=fullflexible,
	frame=single,
	breaklines=true,
	postbreak=\mbox{\textcolor{red}{$\hookrightarrow$}\space},
}
\begin{document}
\title{Experiment `p\_6\_2\_e\_remove\_middle\_array' Results}
\author{\today}
\date{}
\maketitle
\textbf{Experiment outcome: }SUCCESS
\\\textbf{Bad responses: }0
\\\textbf{Responses containing}~\texttt{assume}~\textbf{: }0
\\\textbf{Responses containing}~\texttt{expect}~\textbf{: }0
\\\textbf{Resolution attempts: }2
\\\textbf{Hard fails (resolution): }0
\\\textbf{Soft fails (resolution): }0
\\\textbf{Verification attempts: }2
\section*{Problem Specification}
\textbf{Problem name: }p\_6\_2\_e\_remove\_middle\_array
\\\textbf{Natural language statement: }Write a method to remove the middle element of an array if the array length is odd, or the middle two elements if the length is even.
\\\textbf{Method signature: }p\_6\_2\_e\_remove\_middle\_array(inputs: array<int>) returns (result: array<int>)
\subsection*{Ensures}
\begin{itemize}
\item \texttt{result.Length == if inputs.Length \% 2 == 1 then inputs.Length - 1 else inputs.Length - 2}
\item \texttt{inputs.Length \% 2 == 1 ==> forall i :: 0 <= i < inputs.Length / 2 ==> result[i] == inputs[i] \&\& forall j :: inputs.Length / 2 <= j < result.Length ==> result[j] == inputs[j + 1]}
\item \texttt{inputs.Length \% 2 == 0 ==> forall i :: 0 <= i < inputs.Length / 2 - 1 ==> result[i] == inputs[i] \&\& forall j :: inputs.Length / 2 - 1 <= j < result.Length ==> result[j] == inputs[j + 2]}
\end{itemize}
\subsection*{Requires}
\begin{itemize}
\item \texttt{inputs.Length >= 1}
\end{itemize}
\clearpage
\section*{GenAI interactions}
Below you will find all interactions between the `user' (program) and the `assistant' (OpenAI).
\subsection*{Program $\rightarrow$ GenAI}
\begin{lstlisting}
You are given the following task to perform in Dafny:

Write a method to remove the middle element of an array if the array length is odd, or the middle two elements if the length is even.

The signature should be:

p_6_2_e_remove_middle_array(inputs: array<int>) returns (result: array<int>)

The method should respect the following contract:

requires inputs.Length >= 1, ensures result.Length == if inputs.Length % 2 == 1 then inputs.Length - 1 else inputs.Length - 2, ensures inputs.Length % 2 == 1 ==> forall i :: 0 <= i < inputs.Length / 2 ==> result[i] == inputs[i] && forall j :: inputs.Length / 2 <= j < result.Length ==> result[j] == inputs[j + 1], ensures inputs.Length % 2 == 0 ==> forall i :: 0 <= i < inputs.Length / 2 - 1 ==> result[i] == inputs[i] && forall j :: inputs.Length / 2 - 1 <= j < result.Length ==> result[j] == inputs[j + 2]

Produce and show only the Dafny body of this method, including the curly braces that surround it. Do not show the signature nor contract. You must not use 'assume' nor 'expect' anywhere in your code.
\end{lstlisting}
\subsection*{GenAI $\rightarrow$ Program}
\textbf{System fingerprint: }null
\\\textbf{ID: }chatcmpl-Dn0KpC5ByPqUurB7vQf6RrEhGO7Lu
\\\textbf{Model: }gpt-5.5-2026-04-23
\\\textbf{Created at: }1780571731
\\\textbf{Finish reason: }stop
\\\textbf{Usage: }580 tokens in, and 5475 tokens out
\begin{lstlisting}
{
  var n := inputs.Length;
  if n % 2 == 1 {
    var mid := n / 2;
    assert n == (n / 2) * 2 + n % 2;
    assert n == mid * 2 + 1;
    assert 0 <= mid;
    assert mid <= n - 1;

    result := new int[n - 1];

    var i := 0;
    while i < mid
      invariant 0 <= i <= mid
      invariant result.Length == n - 1
      invariant n == inputs.Length
      invariant forall k :: 0 <= k < i ==> result[k] == inputs[k]
    {
      result[i] := inputs[i];
      i := i + 1;
    }

    var j := mid;
    while j < result.Length
      invariant mid <= result.Length
      invariant mid <= j <= result.Length
      invariant result.Length == n - 1
      invariant n == inputs.Length
      invariant forall k :: 0 <= k < mid ==> result[k] == inputs[k]
      invariant forall k :: mid <= k < j ==> result[k] == inputs[k + 1]
    {
      result[j] := inputs[j + 1];
      j := j + 1;
    }

    assert mid == inputs.Length / 2;
    assert forall k :: 0 <= k < inputs.Length / 2 ==> result[k] == inputs[k];
    assert forall k :: inputs.Length / 2 <= k < result.Length ==> result[k] == inputs[k + 1];
  } else {
    assert 0 <= n % 2 < 2;
    assert n % 2 == 0;
    if n == 1 {
      assert n % 2 == 1;
      assert false;
    }
    assert n >= 2;

    var mid := n / 2;
    assert n == (n / 2) * 2 + n % 2;
    assert n == mid * 2;
    assert 1 <= mid;

    var split := mid - 1;
    assert 0 <= split;
    assert split <= n - 2;

    result := new int[n - 2];

    var i := 0;
    while i < split
      invariant 0 <= i <= split
      invariant result.Length == n - 2
      invariant n == inputs.Length
      invariant forall k :: 0 <= k < i ==> result[k] == inputs[k]
    {
      result[i] := inputs[i];
      i := i + 1;
    }

    var j := split;
    while j < result.Length
      invariant split <= result.Length
      invariant split <= j <= result.Length
      invariant result.Length == n - 2
      invariant n == inputs.Length
      invariant forall k :: 0 <= k < split ==> result[k] == inputs[k]
      invariant forall k :: split <= k < j ==> result[k] == inputs[k + 2]
    {
      result[j] := inputs[j + 2];
      j := j + 1;
    }

    assert split == inputs.Length / 2 - 1;
    assert forall k :: 0 <= k < inputs.Length / 2 - 1 ==> result[k] == inputs[k];
    assert forall k :: inputs.Length / 2 - 1 <= k < result.Length ==> result[k] == inputs[k + 2];
  }
}
\end{lstlisting}
\subsection*{Program $\rightarrow$ GenAI}
\begin{lstlisting}
When using dafny verify, the below error is emitted and verify fails:

problem_data/p_6_2_e_remove_middle_array/program.dfy(13,56): Error: index out of range
   |
13 |     result := new int[n - 1](k => if k < mid then inputs[k] else inputs[k + 1]);
   |                                                         ^

problem_data/p_6_2_e_remove_middle_array/program.dfy(13,56): Error: insufficient reads clause to read array element; Consider extracting inputs[k] to a local variable before the lambda expression, or adding 'reads inputs' in the enclosing lambda specification for resolution
   |
13 |     result := new int[n - 1](k => if k < mid then inputs[k] else inputs[k + 1]);
   |                                                         ^

problem_data/p_6_2_e_remove_middle_array/program.dfy(13,71): Error: index out of range
   |
13 |     result := new int[n - 1](k => if k < mid then inputs[k] else inputs[k + 1]);
   |                                                                        ^

problem_data/p_6_2_e_remove_middle_array/program.dfy(13,71): Error: insufficient reads clause to read array element; Consider extracting inputs[k + 1] to a local variable before the lambda expression, or adding 'reads inputs' in the enclosing lambda specification for resolution
   |
13 |     result := new int[n - 1](k => if k < mid then inputs[k] else inputs[k + 1]);
   |                                                                        ^

problem_data/p_6_2_e_remove_middle_array/program.dfy(53,57): Error: index out of range
   |
53 |     result := new int[n - 2](k => if k < left then inputs[k] else inputs[k + 2]);
   |                                                          ^


Dafny program verifier finished with 1 verified, 5 errors

Correct the error in code by altering only the method body. Produce and show only the Dafny body, including the curly braces that surround it. Do not show the signature nor contract. You must not use 'assume' nor 'expect' anywhere in your code.
\end{lstlisting}
\subsection*{GenAI $\rightarrow$ Program}
\textbf{System fingerprint: }null
\\\textbf{ID: }chatcmpl-Dn0NkK1lLwU54zpC4IRrwWdzyTir3
\\\textbf{Model: }gpt-5.5-2026-04-23
\\\textbf{Created at: }1780571912
\\\textbf{Finish reason: }stop
\\\textbf{Usage: }2046 tokens in, and 6903 tokens out
\begin{lstlisting}
{
  var n := inputs.Length;

  if n % 2 == 1 {
    var mid := n / 2;
    assert 0 <= n % 2 < 2;
    assert n == (n / 2) * 2 + n % 2;
    assert n == mid * 2 + 1;
    assert 0 <= mid;
    assert mid < n;

    result := new int[n - 1];
    assert result != inputs;
    assert result.Length == n - 1;
    assert result.Length == mid * 2;
    assert mid <= result.Length;

    var i := 0;
    while i < mid
      invariant 0 <= i <= mid
      invariant n == inputs.Length
      invariant result != inputs
      invariant result.Length == n - 1
      invariant mid <= result.Length
      invariant mid < n
      invariant forall k :: 0 <= k < i ==> result[k] == inputs[k]
    {
      assert 0 <= i < result.Length;
      assert 0 <= i < inputs.Length;
      result[i] := inputs[i];
      i := i + 1;
    }

    var j := mid;
    while j < result.Length
      invariant mid <= j <= result.Length
      invariant n == inputs.Length
      invariant result != inputs
      invariant result.Length == n - 1
      invariant mid <= result.Length
      invariant forall k :: 0 <= k < mid ==> result[k] == inputs[k]
      invariant forall k :: mid <= k < j ==> result[k] == inputs[k + 1]
    {
      assert 0 <= j < result.Length;
      assert 0 <= j + 1 < inputs.Length;
      result[j] := inputs[j + 1];
      j := j + 1;
    }

    assert mid == inputs.Length / 2;
    assert result.Length == inputs.Length - 1;
    assert result.Length == (if inputs.Length % 2 == 1 then inputs.Length - 1 else inputs.Length - 2);
    assert forall k :: 0 <= k < inputs.Length / 2 ==> result[k] == inputs[k];
    assert forall k :: inputs.Length / 2 <= k < result.Length ==> result[k] == inputs[k + 1];
    assert forall k :: 0 <= k < inputs.Length / 2 ==> result[k] == inputs[k] && (forall t :: inputs.Length / 2 <= t < result.Length ==> result[t] == inputs[t + 1]);
  } else {
    assert n % 2 != 1;
    assert 0 <= n % 2 < 2;
    assert n % 2 == 0;
    if n == 1 {
      assert n % 2 == 1;
      assert false;
    }
    assert n >= 2;

    var mid := n / 2;
    assert n == (n / 2) * 2 + n % 2;
    assert n == mid * 2;
    assert 1 <= mid;

    var left := mid - 1;
    assert 0 <= left;
    assert mid == left + 1;

    result := new int[n - 2];
    assert result != inputs;
    assert result.Length == n - 2;
    assert result.Length == left * 2;
    assert left <= result.Length;

    var i := 0;
    while i < left
      invariant 0 <= i <= left
      invariant n == inputs.Length
      invariant result != inputs
      invariant result.Length == n - 2
      invariant left <= result.Length
      invariant forall k :: 0 <= k < i ==> result[k] == inputs[k]
    {
      assert 0 <= i < result.Length;
      assert 0 <= i < inputs.Length;
      result[i] := inputs[i];
      i := i + 1;
    }

    var j := left;
    while j < result.Length
      invariant left <= j <= result.Length
      invariant n == inputs.Length
      invariant result != inputs
      invariant result.Length == n - 2
      invariant left <= result.Length
      invariant forall k :: 0 <= k < left ==> result[k] == inputs[k]
      invariant forall k :: left <= k < j ==> result[k] == inputs[k + 2]
    {
      assert 0 <= j < result.Length;
      assert 0 <= j + 2 < inputs.Length;
      result[j] := inputs[j + 2];
      j := j + 1;
    }

    assert left == inputs.Length / 2 - 1;
    assert result.Length == inputs.Length - 2;
    assert result.Length == (if inputs.Length % 2 == 1 then inputs.Length - 1 else inputs.Length - 2);
    assert forall k :: 0 <= k < inputs.Length / 2 - 1 ==> result[k] == inputs[k];
    assert forall k :: inputs.Length / 2 - 1 <= k < result.Length ==> result[k] == inputs[k + 2];
    assert forall k :: 0 <= k < inputs.Length / 2 - 1 ==> result[k] == inputs[k] && (forall t :: inputs.Length / 2 - 1 <= t < result.Length ==> result[t] == inputs[t + 2]);
  }
}
\end{lstlisting}
\section*{Final Program}
\begin{lstlisting}
method p_6_2_e_remove_middle_array(inputs: array<int>) returns (result: array<int>)
	requires inputs.Length >= 1
	ensures result.Length == if inputs.Length % 2 == 1 then inputs.Length - 1 else inputs.Length - 2
	ensures inputs.Length % 2 == 1 ==> forall i :: 0 <= i < inputs.Length / 2 ==> result[i] == inputs[i] && forall j :: inputs.Length / 2 <= j < result.Length ==> result[j] == inputs[j + 1]
	ensures inputs.Length % 2 == 0 ==> forall i :: 0 <= i < inputs.Length / 2 - 1 ==> result[i] == inputs[i] && forall j :: inputs.Length / 2 - 1 <= j < result.Length ==> result[j] == inputs[j + 2]
{
  var n := inputs.Length;

  if n % 2 == 1 {
    var mid := n / 2;
    assert 0 <= n % 2 < 2;
    assert n == (n / 2) * 2 + n % 2;
    assert n == mid * 2 + 1;
    assert 0 <= mid;
    assert mid < n;

    result := new int[n - 1];
    assert result != inputs;
    assert result.Length == n - 1;
    assert result.Length == mid * 2;
    assert mid <= result.Length;

    var i := 0;
    while i < mid
      invariant 0 <= i <= mid
      invariant n == inputs.Length
      invariant result != inputs
      invariant result.Length == n - 1
      invariant mid <= result.Length
      invariant mid < n
      invariant forall k :: 0 <= k < i ==> result[k] == inputs[k]
    {
      assert 0 <= i < result.Length;
      assert 0 <= i < inputs.Length;
      result[i] := inputs[i];
      i := i + 1;
    }

    var j := mid;
    while j < result.Length
      invariant mid <= j <= result.Length
      invariant n == inputs.Length
      invariant result != inputs
      invariant result.Length == n - 1
      invariant mid <= result.Length
      invariant forall k :: 0 <= k < mid ==> result[k] == inputs[k]
      invariant forall k :: mid <= k < j ==> result[k] == inputs[k + 1]
    {
      assert 0 <= j < result.Length;
      assert 0 <= j + 1 < inputs.Length;
      result[j] := inputs[j + 1];
      j := j + 1;
    }

    assert mid == inputs.Length / 2;
    assert result.Length == inputs.Length - 1;
    assert result.Length == (if inputs.Length % 2 == 1 then inputs.Length - 1 else inputs.Length - 2);
    assert forall k :: 0 <= k < inputs.Length / 2 ==> result[k] == inputs[k];
    assert forall k :: inputs.Length / 2 <= k < result.Length ==> result[k] == inputs[k + 1];
    assert forall k :: 0 <= k < inputs.Length / 2 ==> result[k] == inputs[k] && (forall t :: inputs.Length / 2 <= t < result.Length ==> result[t] == inputs[t + 1]);
  } else {
    assert n % 2 != 1;
    assert 0 <= n % 2 < 2;
    assert n % 2 == 0;
    if n == 1 {
      assert n % 2 == 1;
      assert false;
    }
    assert n >= 2;

    var mid := n / 2;
    assert n == (n / 2) * 2 + n % 2;
    assert n == mid * 2;
    assert 1 <= mid;

    var left := mid - 1;
    assert 0 <= left;
    assert mid == left + 1;

    result := new int[n - 2];
    assert result != inputs;
    assert result.Length == n - 2;
    assert result.Length == left * 2;
    assert left <= result.Length;

    var i := 0;
    while i < left
      invariant 0 <= i <= left
      invariant n == inputs.Length
      invariant result != inputs
      invariant result.Length == n - 2
      invariant left <= result.Length
      invariant forall k :: 0 <= k < i ==> result[k] == inputs[k]
    {
      assert 0 <= i < result.Length;
      assert 0 <= i < inputs.Length;
      result[i] := inputs[i];
      i := i + 1;
    }

    var j := left;
    while j < result.Length
      invariant left <= j <= result.Length
      invariant n == inputs.Length
      invariant result != inputs
      invariant result.Length == n - 2
      invariant left <= result.Length
      invariant forall k :: 0 <= k < left ==> result[k] == inputs[k]
      invariant forall k :: left <= k < j ==> result[k] == inputs[k + 2]
    {
      assert 0 <= j < result.Length;
      assert 0 <= j + 2 < inputs.Length;
      result[j] := inputs[j + 2];
      j := j + 1;
    }

    assert left == inputs.Length / 2 - 1;
    assert result.Length == inputs.Length - 2;
    assert result.Length == (if inputs.Length % 2 == 1 then inputs.Length - 1 else inputs.Length - 2);
    assert forall k :: 0 <= k < inputs.Length / 2 - 1 ==> result[k] == inputs[k];
    assert forall k :: inputs.Length / 2 - 1 <= k < result.Length ==> result[k] == inputs[k + 2];
    assert forall k :: 0 <= k < inputs.Length / 2 - 1 ==> result[k] == inputs[k] && (forall t :: inputs.Length / 2 - 1 <= t < result.Length ==> result[t] == inputs[t + 2]);
  }
}
\end{lstlisting}
\section*{Total Token Usage}
\textbf{Input tokens: }2626
\\\textbf{Output tokens: }12378
\\\textbf{Reasoning tokens: }10209
\\\textbf{Sum of `total tokens': }15004
\section*{Experiment Timings}
\textbf{Overall Experiment} started at 1780571730484, ended at 1780572086445, lasting 355961ms (355.96 seconds)
\\\textbf{Iteration \#1} started at 1780571730484, ended at 1780571910111, lasting 179627ms (179.63 seconds)
\\\textbf{Iteration \#2} started at 1780571910111, ended at 1780572086445, lasting 176334ms (176.33 seconds)
\end{document}
